% !TEX root = IE3_expB_report.tex
\documentclass[lualatex,a4paper,10pt]{jlreq}
% Shared LuaLaTeX preamble.
\usepackage{luatexja-fontspec}
\setmainfont{TeX Gyre Termes}
\setsansfont{TeX Gyre Heros}
\setmonofont{Latin Modern Mono}
\setmainjfont[BoldFont=HaranoAjiGothic-Medium]{HaranoAjiMincho-Regular}
\setsansjfont[BoldFont=HaranoAjiGothic-Bold]{HaranoAjiGothic-Medium}
\usepackage[top=22mm,bottom=22mm,left=23mm,right=23mm]{geometry}
\usepackage{amsmath,amssymb,booktabs,array,longtable,tabularx}
\usepackage{graphicx,xcolor,tikz,circuitikz}
\usetikzlibrary{arrows.meta,positioning,calc}
\usepackage{enumitem,fancyhdr,listings,needspace}
\usepackage[unicode,colorlinks=true,linkcolor=labblue,urlcolor=labteal]{hyperref}
\usepackage{bookmark}
\definecolor{labblue}{RGB}{30,70,112}
\definecolor{labteal}{RGB}{0,100,105}
\definecolor{labgray}{gray}{0.95}
\ModifyHeading{section}{font={\large\sffamily\bfseries\color{labblue}}}
\ModifyHeading{subsection}{font={\normalsize\sffamily\bfseries\color{labteal}}}
\setlength{\parindent}{1\zw}
\setlength{\parskip}{3pt}
\setlength{\emergencystretch}{2em}
\renewcommand{\arraystretch}{1.35}
\setlist{itemsep=3pt,topsep=4pt,leftmargin=2\zw}
\newcolumntype{Y}{>{\raggedright\arraybackslash}X}
\newcolumntype{C}[1]{>{\centering\arraybackslash}p{#1}}
\pagestyle{fancy}
\fancyhf{}
\fancyhead[L]{\small\color{labblue} IE3 後期・電子工学実験}
\fancyhead[R]{\small テーマ\LabCode}
\fancyfoot[C]{\thepage}
\setlength{\headheight}{16pt}
\DeclareOldFontCommand{\rm}{\normalfont\rmfamily}{\mathrm}
\lstset{basicstyle=\small\ttfamily,columns=fullflexible,keepspaces=true,
 breaklines=true,frame=single,backgroundcolor=\color{labgray},
 numbers=left,numberstyle=\tiny,showstringspaces=false,aboveskip=8pt,belowskip=8pt}
\newcommand{\blank}{\rule{22mm}{0.3pt}}
\newcommand{\answerline}{\par\noindent\rule{\linewidth}{0.25pt}\par}
\newcommand{\writing}[1]{\par\Needspace{\dimexpr#1+5mm\relax}\noindent%
 \fcolorbox{labblue!35}{white}{\parbox[c][#1][t]{\dimexpr\linewidth-2\fboxsep-2\fboxrule\relax}{\vspace{2mm}\noindent\strut\hfill\par}}\par}
\newcommand{\hint}[1]{\par\smallskip\noindent{\sffamily\bfseries\color{labteal}記述の視点：}#1\par}
\newcommand{\note}[1]{\par\smallskip\noindent\fcolorbox{labblue!45}{labblue!5}{%
 \parbox{\dimexpr\linewidth-2\fboxsep-2\fboxrule\relax}{#1}}\par\smallskip}
\newcommand{\eqerror}{\varepsilon=\frac{x_{\mathrm{meas}}-x_{\mathrm{th}}}{x_{\mathrm{th}}}\times100\ \%}
\newcommand{\LabSource}{徳山工業高等専門学校 情報電子工学科，
 『2026年度 電子工学実験（3年後期）テキスト』，テーマ\LabCode，
 \expandafter\path\expandafter{\SourcePdf}。}
\newcommand{\reporttitle}{
 \noindent{\small\sffamily\color{labteal}実験報告書・記入ガイド}\par
 \noindent{\LARGE\sffamily\bfseries \LabCode\quad\LabTitle}\par\smallskip
 \noindent 氏名：\blank\quad 班：\blank\quad 実験日：\blank\par
 \note{目的・原理・手順を確認し、実測値と自分の考察を記入する。
 考察のガイドは、検討する事象・使う測定結果・記述の方針を示す。}
}
\newcommand{\prelabtitle}{
 \noindent{\small\sffamily\color{labteal}事前学習・前提知識プリント}\par
 \noindent{\LARGE\sffamily\bfseries \LabCode\quad\LabTitle}\par\smallskip
 \noindent 氏名：\blank\quad 班：\blank\quad 予習日：\blank\par
 \note{用語、回路の動き、計算、測定表への記入の順に読む。
 例題の値は説明用であり、実験結果ではない。}
}
\newcommand{\tocrow}[3]{\hyperref[#1]{#2}& #3 &\pageref*{#1}\\[2mm]}
\newcommand{\documentmap}[1]{
 \section*{目次と概要}
 \noindent 青色の項目をクリックすると該当箇所へ移動する。\par
 {\small\begin{longtable}{@{}p{0.29\linewidth}p{0.61\linewidth}r@{}}
 \toprule {\color{labblue}\bfseries 項目} & {\color{labblue}\bfseries 読むと分かること} & 頁\\\midrule
 \endhead #1\bottomrule
 \end{longtable}}
 \clearpage
}
\newenvironment{terms}{\par\smallskip\begin{longtable}{@{}p{0.23\linewidth}p{0.73\linewidth}@{}}
 \toprule {\color{labteal}\bfseries 用語}&{\color{labteal}\bfseries 意味・回路での役割}\\\midrule\endhead}
 {\bottomrule\end{longtable}}
\newcommand{\term}[2]{\textbf{#1}&#2\\[2mm]}
\newcommand{\guidepoint}[4]{
 \Needspace{32mm}\subsection{#1}
 \noindent{\bfseries\color{labteal}考える事象：}#2\par
 \noindent{\bfseries\color{labteal}参照する結果：}#3\par
 \noindent{\bfseries\color{labteal}記述の方針：}#4\par
}
\newcommand{\reportclosing}[1]{
 \clearpage\section{結論のまとめ方}\label{sec:closing}
 \hint{#1}
 \ConclusionGuide
 \begin{enumerate}
 \item 目的に対応する最も重要な実測結果を、測定条件と数値を添えて述べる。
 \item 原理と一致した点・一致しなかった点を示す。原因は確認済みの事実と推測を分ける。
 \item 実施範囲で言えることと、未確認の条件を一文ずつまとめる。
 \end{enumerate}
 \note{書き出しの型：「○○の条件で△△を測定したところ、□□となった。
 これは◇◇と整合する。ただし××は確認しておらず、一般化には追加測定が必要である。」
 空欄は自分の実測値と根拠で埋める。}
 \noindent 結論（実験後に記入）：\writing{30mm}
 \noindent 改善案・次に確認したい条件：\writing{16mm}
 \section*{参考文献}\label{sec:references}
 \begin{enumerate}
 \item \LabSource
 \item 使用したデータシート・書籍・ウェブ資料（著者、題名、該当箇所、URL、参照日）を追加する。
 \end{enumerate}
}
\newcommand{\prelabclosing}{
 \section*{準備・出典}\label{sec:prelabsource}
 \begin{itemize}[nosep]
 \item 資料、筆記具、記録用紙、電卓と、実験条件・機器設定の記録欄。
 \item 確認問題の解答と、分からない用語・回路点への具体的な質問。
 \end{itemize}
 {\small\noindent\LabSource\par}
}
\newcommand{\simpleflow}[1]{
 \begin{center}
 \begin{tikzpicture}[node distance=5mm and 5mm,
 every node/.style={font=\small},box/.style={draw=labblue,fill=labblue!4,rounded corners=1mm,
 align=center,minimum height=10mm,text width=30mm}]
 #1
 \end{tikzpicture}
 \end{center}
}

\newcommand{\LabCode}{B}
\newcommand{\LabTitle}{共振回路}
\newcommand{\SourcePdf}{IE3_expB_A4.pdf}
\newcommand{\ConclusionGuide}{直列と並列について$f_0,\Delta f,Q_{\rm BW}$を条件とともに並べる。どちらの量がピークだったか、理論との差はどの程度かを結論に含める。
鋭さの比較は同じ測定・負荷条件に基づいて行い、素子誤差や損失を測っていない場合は原因を確定しない。\par}
\begin{document}
\reporttitle
\documentmap{
\tocrow{sec:auto1}{1 目的}{LCの周波数応答と共振の鋭さを実験で確かめる。}
\tocrow{sec:auto2}{2 原理}{LCの周波数応答と共振の鋭さの仕組みと成立条件を整理する。}
\tocrow{sec:auto3}{3 装置・条件}{機器・定数・測定点を確認し、資料の順序で操作する。}
\tocrow{sec:auto4}{4 方法}{機器・定数・測定点を確認し、資料の順序で操作する。}
\tocrow{sec:auto5}{5 測定結果}{周波数曲線、半値点、幅とQを表や図に記録する。}
\tocrow{sec:auto6}{6 考察：結果を根拠に説明する}{周波数曲線、半値点、幅とQを根拠に、比較と原因の検討を進める。}
\tocrow{sec:closing}{7 結論のまとめ方}{主要な結果、成立範囲、未確認の条件を短くまとめる。}
\tocrow{sec:references}{参考文献}{使用した資料と外部の根拠を示す。}
}

\section{目的}\label{sec:auto1}
直列および並列のRLC回路の周波数特性を測定し、共振周波数、共振の鋭さ、半値幅の関係を確かめる。
\section{原理}\label{sec:auto2}
直列回路のインピーダンスと共振条件は
\[
 Z=r+j\left(\omega L-\frac1{\omega C}\right),\qquad
 f_0=\frac1{2\pi\sqrt{LC}},\qquad Q=\frac{\omega_0L}{r}
      =\frac{f_0}{f_2-f_1}
\]
である。$r$にはコイルの損失など直列損失を含む。電流が最大値の$1/\sqrt2$になる周波数を$f_1,f_2$とする。
実際のコイルを$r+j\omega L$で表した並列回路では
\[
 Y=j\omega C+\frac1{r+j\omega L},\qquad
 \omega_{0p}^2=\frac1{LC}-\frac{r^2}{L^2}.
\]
高$Q$では並列共振時のインピーダンスはおよそ$\omega_0^2L^2/r$となる。
\section{装置・条件}\label{sec:auto3}
\par\noindent\begin{tabularx}{\linewidth}{lYY}
\toprule & 直列共振 & 並列共振\\\midrule
$L,C$ & $50\,\mathrm{mH},\ 0.02\,\mu\mathrm{F}$ & 同左\\
測定用抵抗$R$ & $5\,\Omega$ & $5\,\mathrm{k}\Omega$\\
発振器出力 & 実際の設定を記録 & $5\,\mathrm{V_{pp}}$\\
\bottomrule
\end{tabularx}
発振器、オシロスコープ、周波数計、部品の型番・実測値・使用レンジ：
\writing{14mm}
\clearpage
\section{方法}\label{sec:auto4}
\begin{enumerate}
\item 資料の測定回路を配線し、共振回路端子電圧$e_a$と測定抵抗電圧$e_b$を観測する。接続図とプローブの位置を記録する。
\item 周波数を変え、共振点付近では間隔を細かくする。出力振幅を確認しながら$e_a,e_b$と周波数を記録する。
\item $|Z|=R|e_a|/|e_b|$、$|Y|=|e_b|/(R|e_a|)$を求める。双方の電圧は同じ定義（実効値または振幅）で測定する。
\item 直列では$|Y|$、並列では$|Z|$を縦軸、周波数を対数横軸とする。ピークの$1/\sqrt2$の点から半値幅を求める。
\end{enumerate}
\section{測定結果}\label{sec:auto5}
表1：直列共振（不足行は追加する）。
\par\noindent\begin{tabularx}{\linewidth}{YYYYY}
\toprule $f$ / Hz & $e_a$ / V & $e_b$ / V & $|Z|$ / $\Omega$ & $|Y|$ / S\\\midrule
 & & & &\\ & & & &\\ & & & &\\ & & & &\\ & & & &\\ & & & &\\\bottomrule
\end{tabularx}
表2：並列共振。
\par\noindent\begin{tabularx}{\linewidth}{YYYYY}
\toprule $f$ / Hz & $e_a$ / V & $e_b$ / V & $|Z|$ / $\Omega$ & $|Y|$ / S\\\midrule
 & & & &\\ & & & &\\ & & & &\\ & & & &\\ & & & &\\ & & & &\\\bottomrule
\end{tabularx}
\clearpage
図1：直列共振のアドミタンス曲線（測定点と半値点）。
\writing{40mm}
図2：並列共振のインピーダンス曲線（測定点と半値点）。
\writing{40mm}
表3：共振特性の整理。
\par\noindent\begin{tabularx}{\linewidth}{YYYYY}
\toprule 回路 & $f_0$/Hz & $f_1,f_2$/Hz & 半値幅/Hz & $Q$\\\midrule
直列 & & & &\\並列 & & & &\\\bottomrule
\end{tabularx}
\clearpage\section{考察：結果を根拠に説明する}\label{sec:auto6}
\guidepoint{共振周波数の理論との比較}
{直列・並列のピーク周波数が、理想LCの計算値からどの程度ずれるか。}
{表1・2の測定点、図1・2のピーク、表3の$f_0$、実際のLとC。}
{理論値と実測の定義を示して相対差を求める。並列では実コイル損失のある共振条件と理想条件の違いも検討する。Z最大点と虚部零の条件が完全に同じとは限らないことに注意する。}
\guidepoint{二つの方法で求めるQ}
{曲線の幅から得るQと、損失抵抗から見積もるQが一致するか。}
{図の半値点、表3の幅、コイル抵抗測定値、発振器・プローブの条件。}
{まず同じ回路モデルで比較する。直列損失rと測定用Rを区別し、外部負荷を含む幅であるか説明する。並列の幅によるQは回路・高Q近似の適用範囲も明記する。}
\guidepoint{周波数間隔と半値点の読み取り}
{測定点の間隔が広いことでピークや幅を取り逃がしていないか。}
{半値点前後の実測周波数・振幅、グラフの対数目盛、再測定点。}
{補間した点と実測点を区別する。読める周波数の幅を示し、Qがその範囲でどの程度変わるかを見積もる。追加測定していなければ、その必要性を改善案にする。}
\guidepoint{ピーク量と損失の物理的意味}
{直列でY、並列でZがピークになる理由と、有限の高さとなる理由。}
{二曲線の形、ピーク値、損失・負荷条件、端子電圧の変化。}
{「流れやすい／流れにくい」を測る量に結び付けて説明する。直列の半値点は一定端子電圧・一定rなら損失電力が半分。並列では同じ電力説明を無条件に当てはめず、測ったZの指標として述べる。}
\subsection{資料の課題・追加検討}
\begin{enumerate}
\item 半値幅から求めた$Q$と損失抵抗から求めた$Q$を比較し、差の原因を検討する。
\item 直列共振で電流が$1/\sqrt2$となる条件から$\Delta f=f_0/Q$を導く。
\item $L=800\,\mathrm{mH},C=200\,\mathrm{pF},Q=50$の共振周波数と半値幅を計算する。結合回路の付録を扱った場合は結合度と曲線形状を整理する。
\end{enumerate}
\writing{25mm}
\reportclosing{直列と並列でピークとなる量を明記し、共振周波数と$Q$を数値で要約する。}
\end{document}
